\documentclass[a4paper]{article}
\usepackage[pdftex]{graphicx}
\usepackage[utf8]{inputenc}
\usepackage{enumerate}
\usepackage{siunitx}
\sisetup{locale=DE} 
\usepackage{href-ul}
\hypersetup{
	colorlinks=true,
	linkcolor=blue,
	urlcolor=blue}
\usepackage{icomma}
\usepackage{amssymb}
\usepackage{geometry}
\geometry{a4paper, top=15mm, left=15mm, right=15mm, bottom=15mm,
	headsep=10mm, footskip=12mm}

%Source: img_0236

% Start the document
\begin{document}
%Titel
{\bf Mechanische Wellen }\\
\begin{enumerate}[1.]
%Aufgabe 1
\item Durch einen harmonisch schwingenden Erreger wird in einer Wellenwanne eine Wasserwelle mit geraden Wellenfronten erzeugt. Die Wellenbewegung hat bereits alle Punkte der Wasseroberfläche erfasst. Die x-Achse eines kartesischen Koordinatensystems verläuft horizontal und senkrecht zu den Wellenfronten, die y-Achse des Koordinatensystems ist vertikal nach oben gerichtet. Der Erreger der Wasserwelle befindet sich an der Stelle $x_0=\SI{0}{\centi\meter}$ und bewegt sich zum Zeitpunkt $t_0=0s$ gerade durch die Nulllage nach oben. Die Welle, die sich vom Erreger aus in positiver x-Richtung ausbreitet, lässt sich durch folgende Wellengleichung beschreiben: 
$$y(t,x)=\SI{0,70}{\centi\meter} \cdot \sin{ \left[ 2 \pi \left( {t \over \SI{0,40}{\second}}-{x \over \SI{6,0}{\centi\meter}} \right) \right]}$$

\begin{enumerate}[a)]
\item Berechnen Sie die Frequenz $f$ der Erregerschwingung und den Betrag $c$ der Ausbreitungsgeschwindigkeit der Welle aus den Daten der Wellen\-gleichung.
\item Zeichnen Sie im Maßstab 1:1 für den Zeitpunkt $t_1=\SI{0,30}{\second}$ das Momentanbild der Welle im
Bereich $ \SI{0}{\centi\meter} \le x \le \SI{12,0}{\centi\meter}$
\item $P$ sei ein Punkt an der Wasseroberfläche mit der x-Koordinate $x_P=\SI{10,0}{\centi\meter}$. Berechnen Sie für den Zeitpunkt $t_1=\SI{0,30}{\second}$ den Betrag der Geschwindigkeit des Punktes $P$. Zeichnen Sie die Geschwindigkeit $\vec{v}~ (\SI{0,30}{\second}; \SI{10,0}{\centi\meter})$ in das Momentanbild von Teilaufgabe 1. b) ein, indem Sie den zugehörigen Vektorpfeil am Punkt $P$ anheften. Verwenden Sie dabei den Maßstab 
$\SI{5,0}{\centi\meter\second^{-1}} ~\hat{=} ~\SI{1}{\centi\meter}$
\item Zeichnen Sie auch den Vektorpfeil für die Geschwindigkeit $\vec{v} ~(\SI{0,30}{\second}; \SI{1,5}{\centi\meter})$ unter Verwendung des unter 1. c) angegebenen Maßstabes in das Bild von Teilaufgabe 1. b) ein.
\end{enumerate}

%Aufgabe 2
\item Die Welle mit geraden Wellenfronten aus 1. trifft nun auf eine senkrecht zur Ausbreitungsrichtung der Welle stehende Wand mit zwei engen Spalten $S_1$ und $S_2$. Der Abstand der Spaltmitten beträgt $d=\SI{20,0}{\centi\meter}$. Die von den beiden Spalten ausgehenden Wellen breiten sich mit einer Geschwindigkeit vom Betrag $c=\SI{15}{\centi\meter\second^{-1}}$ aus. Zur Vereinfachung soll angenommen werden, dass jede der beiden Wellen für sich allein jeden beliebigen Punkt der Wasseroberfläche hinter der Wand zu einer Schwingung mit der Amplitude $A=\SI{0,70}{\centi\meter}$ anregt. 
  
\begin{enumerate}[a)]
\item Ein Punkt $Q$ der Wasseroberfläche hat vom Spalt $S_1$ den Abstand $\SI{16,0}{\centi\meter}$ und von Spalt $S_2$ den Abstand $\SI{28,0}{\centi\meter}$. Die von den beiden Spalten ausgehenden Wellen haben bereits den Punkt $Q$ erfasst. Beschreiben Sie die Bewegung des Punktes $Q$. Erklären Sie, wie die beschriebene Erscheinung zu Stande kommt.
\item Wenn man die Frequenz der Erregerschwingung vom unter 1.a) berechneten Wert ausgehend steigert, bleibt der Punkt $Q$ bei bestimmten Frequenzen praktisch in Ruhe. Berechnen Sie die Frequenz $f_1$, bei der dies das erste Mal der Fall ist.
\item Die Anzahl der Kurven mit maximaler Schwingungsamplitude ist abhängig von der Frequenz der Erregerschwingung. Bestimmen Sie den Frequenzbereich, für den genau sieben Kurven mit maximaler Schwingungsamplitude beobachtet werden können.  
\end{enumerate} 

\end{enumerate} 
\fbox{
	\begin{minipage}{0.5\textwidth}
		Weitere Arbeitsblätter gibt es unter 
		
		\href{https://www.okuyakl.de}{www.okuyakl.de}
	\end{minipage}
	\hfill
	\begin{minipage}{0.4\textwidth}
		\includegraphics[width=1.5 cm]{../../viecher/zwe03}
		\includegraphics[width=2 cm]{../../viecher/afanticon1}
		
\end{minipage}}
%Ende
\end{document}